\subsection{Agent-Based Cell--ECM Interaction Model}
\label{sec:ecm_confinement}

The extracellular matrix (ECM) surrounding an organoid provides both structural confinement and adhesive anchorage.
In the present model, the ECM is discretised onto a spatial grid and each grid element is treated as an independent mechanical agent that interacts with nearby cells through linear springs.
This agent-based formulation captures two essential properties of the basement membrane that a simple radial boundary condition cannot: tangential adhesion, which anchors the tissue in place, and density-dependent force transmission, which couples matrix degradation to local mechanical release.

\subsubsection{ECM Field Representation}

The ECM is stored on a two-dimensional grid (square or hexagonal) with spacing $h$ that tiles the simulation domain.
Each grid element $j$ carries a scalar density field $\rho_j \in [0,1]$, representing the local concentration of matrix material, and a fibre orientation angle $\theta_j$ for optional contact guidance.
The grid is initialised with $\rho_j = 1$ outside the organoid boundary and $\rho_j = 0$ inside the lumen and epithelial layer, reflecting the biological arrangement in which basement membrane components such as laminin and collagen~IV are deposited basolaterally \cite{Yurchenco2011}.

\subsubsection{Cell--ECM Spring Force Law}

Each cell $i$ located at position $\mathbf{x}_i$ interacts with all ECM elements whose centres $\mathbf{x}_j^{\text{ECM}}$ lie within a cutoff distance $d_{\text{cut}}$.
For each such pair, the force on cell $i$ due to ECM element $j$ is given by

\begin{equation}
    \mathbf{F}_{ij}^{\text{ECM}} = k_{\text{ECM}}\;\rho_j\;\bigl(d_{ij} - s_0\bigr)\;\hat{\mathbf{d}}_{ij},
    \label{eq:ecm_spring}
\end{equation}

\noindent where
\begin{itemize}
    \item $k_{\text{ECM}}$ is the cell--ECM spring stiffness,
    \item $\rho_j \in [0,1]$ is the ECM density at element $j$, acting as a force multiplier,
    \item $d_{ij} = \lVert \mathbf{x}_j^{\text{ECM}} - \mathbf{x}_i \rVert$ is the distance between the cell and the ECM element,
    \item $s_0$ is the spring rest length,
    \item $\hat{\mathbf{d}}_{ij} = (\mathbf{x}_j^{\text{ECM}} - \mathbf{x}_i) / d_{ij}$ is the unit vector from cell to ECM element.
\end{itemize}

\noindent The total ECM force on cell $i$ is the sum over all interacting elements:

\begin{equation}
    \mathbf{F}_i^{\text{ECM}} = \sum_{j:\, d_{ij} < d_{\text{cut}}} k_{\text{ECM}}\;\rho_j\;\bigl(d_{ij} - s_0\bigr)\;\hat{\mathbf{d}}_{ij}.
    \label{eq:ecm_total}
\end{equation}

A key design decision is that the force law in Equation~\eqref{eq:ecm_spring} is \emph{purely linear} in the overlap $(d_{ij} - s_0)$.
Standard cell--cell spring forces in node-based models employ a logarithmic repulsion term for overlapping cells, $F \propto \ln(1 + \delta/s_0)$, which diverges as $\delta \to -s_0$, and an exponentially decaying attraction, $F \propto \delta\,\exp(-\alpha\,\delta/s_0)$, for separated cells \cite{Meineke2001}.
These nonlinearities are appropriate for cell--cell contact, where volume exclusion creates a hard-core repulsion.
However, for cell--ECM interaction, the biological reality is that cells migrate \emph{through} the matrix by proteolytic degradation of ECM fibres via matrix metalloproteinases (MMPs) and by squeezing through pores in the network \cite{Wolf2013}.
A divergent repulsive force would artificially prevent this behaviour.
The linear law permits cells to overlap with ECM elements at a finite, bounded force, with the physical resistance controlled instead by the density-dependent prefactor $\rho_j$: degraded matrix ($\rho_j \to 0$) offers negligible resistance regardless of overlap.

\subsubsection{Density-Dependent Force Scaling}

The inclusion of $\rho_j$ as a multiplicative factor in Equation~\eqref{eq:ecm_spring} creates a natural coupling between ECM remodelling and mechanical confinement.
In regions where matrix metalloproteinase (MMP) secretion has reduced the local density, the spring force vanishes proportionally.
This produces a mechanochemical feedback loop: proliferating cells crowd toward the ECM boundary, secrete MMPs that locally degrade the matrix, and the resulting reduction in confinement force permits outward budding.
The density evolves via

\begin{equation}
    \rho_j^{n+1} = \max\!\bigl(0,\;\rho_j^n - k_{\text{deg}}\,\Delta t\bigr),
    \label{eq:ecm_degradation}
\end{equation}

\noindent where $k_{\text{deg}}$ is the degradation rate per cell per unit time and the $\max$ operator prevents unphysical negative densities.
A diffusion step smooths the density field at each timestep, representing cross-linking and lateral spreading of matrix material \cite{Lu2011}:

\begin{equation}
    \rho_j^{n+1} \leftarrow \rho_j^{n+1} + D\,\Delta t\,\bigl(\bar{\rho}_j - \rho_j^{n+1}\bigr),
    \label{eq:ecm_diffusion}
\end{equation}

\noindent where $D$ is the diffusion coefficient and $\bar{\rho}_j$ is the mean density of the neighbouring grid elements.

\subsubsection{Tangential Anchoring and Rotational Stability}

Treating each ECM element as an individual spring agent, rather than applying a single radial force directed toward the tissue centroid, introduces force components in all spatial directions.
Consider a cell sitting at the organoid boundary: it is connected by springs to multiple ECM elements arranged around its exterior.
If the organoid were to undergo rigid-body rotation, each cell would move tangentially relative to its anchoring ECM elements, stretching those springs and generating restoring forces that oppose the rotation.
Similarly, a bulk translation of the entire tissue stretches springs asymmetrically, producing a net restoring force.
This emergent resistance to rigid-body motion arises naturally from the spatial distribution of spring connections and does not require any explicit penalty term.

This property is a direct consequence of modelling the ECM as a discrete viscoelastic material rather than an isotropic pressure.
Experimental studies of Matrigel-embedded organoids confirm that the basement membrane provides both radial confinement and tangential adhesion through integrin--laminin binding \cite{Gjorevski2016}, and that disruption of these adhesions (e.g.\ via integrin-blocking antibodies) permits organoid drift and deformation.
The agent-based spring model captures this dual mechanical role without requiring separate force terms for confinement and adhesion.

\subsubsection{Summary of Parameters}

Table~\ref{tab:ecm_params} lists the configurable parameters for the cell--ECM interaction force.

\begin{table}[h]
    \centering
    \caption{Parameters governing the cell--ECM confinement force.}
    \label{tab:ecm_params}
    \begin{tabular}{lll}
        \hline
        \textbf{Symbol} & \textbf{Parameter} & \textbf{Default} \\
        \hline
        $k_{\text{ECM}}$  & \texttt{ecmConfinementStiffness} & 5.0 \\
        $s_0$             & \texttt{ecmSpringRestLength}     & 0.0 \\
        $d_{\text{cut}}$  & \texttt{ecmInteractionCutoff}    & 1.5 \\
        $h$               & \texttt{ecmGridSpacing}          & 0.25 \\
        $k_{\text{deg}}$  & \texttt{ecmDegradationRate}      & 0.5 \\
        $D$               & \texttt{ecmDiffusionCoeff}       & 0.1 \\
        \hline
    \end{tabular}
\end{table}
